% Generated draft: requires mathematical and source review.
\documentclass[11pt]{article}
\usepackage[margin=25mm]{geometry}
\usepackage{fontspec}
\setmainfont{Latin Modern Roman}
\usepackage{amsmath,amssymb,mathtools,unicode-math}
\setmathfont{Latin Modern Math}
\usepackage{enumitem,needspace,xurl,hyperref}
\setlength{\parindent}{0pt}
\setlength{\parskip}{5pt}
\setlength{\emergencystretch}{3em}
\newcommand{\N}{\mathbb{N}}
\newcommand{\Z}{\mathbb{Z}}
\newcommand{\Q}{\mathbb{Q}}
\newcommand{\R}{\mathbb{R}}
\newcommand{\C}{\mathbb{C}}
\newcommand{\F}{\mathbb{F}}
\newcommand{\A}{\mathbb{A}}
\newcommand{\PP}{\mathbb{P}}
\newcommand{\E}{\mathbb{E}}
\newcommand{\eps}{\varepsilon}
\newcommand{\dd}{\,\mathrm{d}}
\begin{document}
\section*{UnsolvedMath ID: 30000697}
\textbf{Source problem code:} catalogue entry
\subsection{Definitions and mathematical statement}
Let \(\gamma\) be a \(\sigma\)-finite Borel measure on \((0,\infty)\), as stated in the source. In the usual \(d\)-dimensional setting, a L\'evy measure \(\rho\) on \(\R^d\) satisfies
\[
 \rho(\{0\})=0,\qquad
 \int_{\R^d}(1\wedge\|z\|^2)\,\rho(dz)<\infty.
\]
The dimension and underlying space of \(\rho\) are not explicitly supplied, so this standard finite-dimensional convention is made conditional. For a Borel set \(A\), put \(Ax^{-1}=\{z:xz\in A\}\). The source's transform is the positive scaling mixture
\[
 \Upsilon_\gamma(\rho)(A)=\int_0^\infty\rho(Ax^{-1})\,\gamma(dx).
\]
This is inverse scaling of the set, or equivalently forward multiplication of points by \(x\), not the reverse convention.

With \(\mathfrak L_d\) denoting the L\'evy measures above, its natural L\'evy domain is
\[
 \mathcal D_\gamma=\{\rho\in\mathfrak L_d:
 \Upsilon_\gamma(\rho)\in\mathfrak L_d\}.
\]
Under this interpretation the integrability requirement can be written
\[
 \int_{\R^d}\int_0^\infty(1\wedge\|xz\|^2)\,\gamma(dx)\rho(dz)<\infty.
\]
The source asks for general conditions on \(\gamma\) ensuring
\[
 \forall\rho_1,\rho_2\in\mathcal D_\gamma,\qquad
 \Upsilon_\gamma(\rho_1)=\Upsilon_\gamma(\rho_2)
 \ \Longrightarrow\ \rho_1=\rho_2.
\]
Equalities are equalities of measures, not merely of finitely many moments or transform values. The entire intended L\'evy domain is the target; no density, moment or support restriction on its elements is added. Whether a criterion is sought for one fixed dimension or uniformly in all dimensions must be specified. A domain with only the zero measure would make injectivity vacuous, and no nontrivial-domain requirement is printed.

The statement assumes only \(\sigma\)-finiteness of the scaling measure \(\gamma\). The literature assessment speaks of a criterion for every L\'evy measure \(\gamma\), which is a narrower or differently normalized class; no additional L\'evy integrability of \(\gamma\) is silently imposed. Its triage checked 2026-08-21 reports domain descriptions, injectivity reductions, sufficient criteria and noninjective examples, but no complete general criterion. Those claims, the partially-solved label and notation-restoration verification are retained as dated metadata, not independently checked current results. No Mellin/Fourier nonvanishing condition is guessed from the scaling formula alone.
\subsection{Short English statement}
Find conditions on the scaling measure that make the Upsilon mixture injective on all input L\'evy measures for which the output remains a L\'evy measure. The inverse-set scaling is explicit. The source allows general \(\sigma\)-finite scaling measures, despite narrower wording in its assessment, and leaves the ambient dimension unspecified.
\subsection{Sources}
\begin{itemize}
\item Local UnsolvedMath registry, permanent ID 30000697, source code catalogue entry, ``Injectivity Criteria for Upsilon Transforms,'' including the full scaling formula, L\'evy-domain question and verification metadata.
\item \emph{Mini-Workshop: L\'evy Processes and Related Topics in Modelling (2007)}, Oberwolfach Reports, pp.\ 417--446, DOI 10.4171/owr/2007/07: \url{https://doi.org/10.4171/owr/2007/07}.
\item Supplied reference ``general Upsilon transformations'': \url{https://alea.impa.br/articles/v4/04-07.pdf}.
\item Supplied reference ``injectivity criteria and counterexample'': \url{https://data.math.au.dk/publications/thiele/2008/imf-thiele-2008-02.pdf}. These URLs are cited as supplied; no PDFs have been opened.
\item Literature triage checked 2026-08-21; record updated 2026-09-16. No external source has been visited or independently checked here.
\end{itemize}
\end{document}
